\section{Additive 2-matroids of small rank or ground-set size}
\label{sec:two-matroid}

Following the additive augmentation convention of Xie and Xu \cite{XX}, a
\emph{2-matroid} is a nonempty downset $\mathcal I\subseteq2^V$ such that
\begin{equation}\label{eq:two-augmentation}
S,T\in\mathcal I,\quad |T|\ge |S|+2
\quad\Longrightarrow\quad
S\cup\{x\}\in\mathcal I\text{ for some }x\in T\setminus S.
\end{equation}
Its rank is $\max\{|S|:S\in\mathcal I\}$. For element weights $w_v\ge0$, write
\[
A_j=\sum_{\substack{S\in\mathcal I\\|S|=j}}w(S).
\]

\begin{theorem}\label{thm:two-matroid-small}
The weighted rank sequence of a 2-matroid on $n$ ground elements is ultra
log-concave of ambient order $n$ if its rank is at most four or if $n\le7$.
Thus, in either case,
\[
k(n-k)A_k^2\ge(k+1)(n-k+1)A_{k-1}A_{k+1}
\qquad(1\le k\le n-1).
\]
\end{theorem}

We establish the required antipodal-pairs inequalities analytically. On an
even ground set of size $2s$, put
\[
F_j=\{S\in\mathcal I:|S|=j,\ V\setminus S\in\mathcal I\}.
\]
The inequality to prove is $s|F_s|\ge(s+1)|F_{s-1}|$. On two elements this is
automatic for any downset: if $F_0$ is nonempty, the full two-set is feasible,
so both singletons are feasible and $|F_1|=2$.

\begin{lemma}[Four elements]\label{lem:two-matroid-four}
Every 2-matroid on four elements satisfies $2|F_2|\ge3|F_1|$.
\end{lemma}
\begin{proof}
Assume $F_1\ne\varnothing$, as otherwise the assertion is immediate. A
feasible singleton and its feasible complementary triple show that all four
singletons are feasible. Thus $|F_1|=t$, the number of feasible triples.
View the feasible pairs as a graph.

If $t=1$, augment the singleton outside the feasible triple towards that
triple. This gives a feasible pair joining the outside vertex to a vertex
of the triple; its complementary pair is feasible by downward closure.
Hence $|F_2|\ge2$. If $t=2$, the pair shadows of the two triples contain
all edges of $K_4$ except at most one. At least two perfect matchings remain,
so $|F_2|\ge4$. If $t\ge3$, all six pairs are feasible, so $|F_2|=6$.
These three cases give $2|F_2|\ge3t$.
\end{proof}

\begin{lemma}[Six elements]\label{lem:two-matroid-six}
Every 2-matroid on six elements satisfies $3|F_3|\ge4|F_2|$.
\end{lemma}
\begin{proof}
Let $R$ be the simple graph with edge set $F_2$, and put $r=|E(R)|=|F_2|$.
Let $m$ count unordered feasible partitions of $V$ into two triples. Thus
$|F_3|=2m$, and the target is $3m\ge2r$. The case $r=0$ is immediate.

Let $g$ count unordered balanced partitions with an $R$-edge in each part.
Every such partition is feasible: if $e\subseteq A$ and $f\subseteq B$ are
edges, then $A\subseteq V\setminus f$ and $B\subseteq V\setminus e$, and
the complementary four-set of each $R$-edge is feasible. Hence $m\ge g$.

Suppose first that no edge of $R$ is a vertex cover. Then every edge has an
edge in its complement. For the doubly dependent counts of the ordinary
graph $R$, this says $b_2(R)=r$ and $b_3(R)=2g$. Theorem~\ref{thm:apu}
gives $3b_3(R)\ge4b_2(R)$, so $3m\ge3g\ge2r$.

It remains to consider a vertex-cover edge $ab$. The four-set
$L=V\setminus\{a,b\}$ is independent in $R$. Put
\[
A=N_R(a)\cap L,\quad B=N_R(b)\cap L,\quad p=|A|,\quad q=|B|,
\]
and interchange $a,b$ so that $0\le q\le p\le4$. All edges are $ab$ or
join a centre to a marked leaf, so $r=1+p+q$. Since $ab$ and $L$ are
feasible, augmentation gives a feasible triple $abx$ for some $x\in L$.
Its complement is feasible by downward closure. This partition is not
counted by $g$, because its other part lies in the independent set $L$.
Therefore
\begin{equation}\label{eq:two-matroid-extra}
m\ge g+1.
\end{equation}

For $q\ge1$, the six balanced partitions separating the centres are
parametrised by two-sets $X\subseteq L$, giving parts $\{a\}\cup X$ and
$\{b\}\cup(L\setminus X)$. They contribute to $g$ exactly when $X$ meets
$A$ and $L\setminus X$ meets $B$. Consequently
\[
g\ge6-\bin{4-p}{2}-\bin{4-q}{2}.
\]
Here $\bin{4-p}{2}$ counts the two-sets $X$ that miss $A$, and
$\bin{4-q}{2}$ counts those whose complement $L\setminus X$, itself a
two-set, misses $B$; the bound is a union bound over these two failure modes.
If $p=q=1$ and $A\ne B$, direct counting gives $g=2$. The exceptional
case $p=q=1$, $A=B$ is a triangle with isolated vertices and is treated
below. All other cases with $q\ge1$ are exhausted by the following table.
The last column, together with \eqref{eq:two-matroid-extra}, proves the
desired inequality in each case.
\begin{center}
\begin{tabular}{rrrr}
\toprule
$p$ & $q$ & Lower bound $g_0$ for $g$ & $3(g_0+1)-2(1+p+q)$\\
\midrule
1 & 1 & 2 & 3\\
2 & 1 & 2 & 1\\
3 & 1 & 3 & 2\\
4 & 1 & 3 & 0\\
2 & 2 & 4 & 5\\
3 & 2 & 5 & 6\\
4 & 2 & 5 & 4\\
3 & 3 & 6 & 7\\
4 & 3 & 6 & 5\\
4 & 4 & 6 & 3\\
\bottomrule
\end{tabular}
\end{center}

If $q=0$, the graph $R$ is a star with centre $a$, nonempty leaf set $U$,
and possible isolated vertices; here $|U|=r$. On $W=V\setminus\{a\}$,
define a graph $J$ by $ux\in E(J)$ exactly when $\{a,u,x\}$ is feasible.
Each $u\in U$ has positive degree in $J$, by augmenting $au$ towards its
feasible complementary four-set. Moreover, a degree-one vertex $u\in U$
can have its unique neighbour only outside $U$. Indeed, for each $v\in U$,
augmenting $au$ towards $V\setminus\{a,v\}$ supplies a neighbour of $u$
other than $v$.

Distribute one unit of charge from each $u\in U$ equally among its incident
edges in $J$. An edge within $U$ receives at most one, since both endpoints
have degree at least two; an edge with exactly one endpoint in $U$ also
receives at most one. Thus at least $|U|$ edges of $J$ meet $U$. Each such
edge $ux$, with $u\in U$, gives a feasible balanced partition: the triple
$aux$ is feasible, and its complement lies in the feasible four-set
$V\setminus\{a,u\}$. Distinct edges give distinct unordered partitions,
identified by their part containing $a$. Therefore $m\ge|U|=r$.

Finally, suppose $R$ is a triangle on $a,b,c$ with three isolated vertices
$L$. For each triangle edge, say $ab$, augment $ab$ towards the feasible
four-set $V\setminus ac=\{b\}\cup L$. This supplies $x\in L$ with $abx$
feasible. Its complement is a subset of the feasible four-set
$V\setminus ab$. The three triangle edges yield three distinct unordered
partitions: the part with two triangle vertices identifies the edge.
Hence $m\ge3=r$. This completes all cases.
\end{proof}

\begin{proof}[Proof of Theorem~\ref{thm:two-matroid-small}]
Deletion and contraction by a feasible fixed set preserve
\eqref{eq:two-augmentation}: for contraction by $P$, apply the original
axiom to $P\cup S$ and $P\cup T$, whose cardinality difference is unchanged.
Every positive-probability coordinate restriction therefore remains a
2-matroid, and its rank cannot increase.

For rank at most four, an even conditional ground set of size $2s\ge8$
has $F_{s-1}=\varnothing$, because the complementary set would have size
$s+1\ge5$. Its antipodal-pairs inequality is automatic. The two-point
observation and Lemmas~\ref{lem:two-matroid-four}--\ref{lem:two-matroid-six}
cover the remaining even sizes. If the original ground set instead has
size at most seven, every even conditional ground set already has size at
most six, and the same lemmas apply regardless of rank.

For positive weights, the law proportional to $w(S)\ind_{\{S\in\mathcal I\}}$
has antipodal products equal to the common factor $w(V)/Z^2$ on feasible
complementary pairs. The preceding inequalities therefore give CAPP for
this law, including after every coordinate restriction. A nonempty
downset has interval rank support. The Kahn--Neiman criterion \cite{KN}
gives the asserted weighted ultra log-concavity. Nonnegative weights follow
by continuity, with the same ambient normalisation $\bin nk$.
\end{proof}

\begin{corollary}[Initial coefficients in arbitrary rank]
\label{cor:two-matroid-initial}
For any 2-matroid on $n$ ground elements and any nonnegative element
weights, its coefficients satisfy
\[
k(n-k)A_k^2\ge(k+1)(n-k+1)A_{k-1}A_{k+1}
\qquad(1\le k\le\min(3,n-1)).
\]
Thus its first $\min(5,n+1)$ binomially normalised coefficients are
log-concave, using the original ambient order $n$.
\end{corollary}
\begin{proof}
The rank-four truncation $\{S\in\mathcal I:|S|\le4\}$ is a 2-matroid on
the same ground set. An augmentation between its feasible $S,T$ has size
$|S|+1\le|T|-1\le3$, so remains in the truncation. Its coefficients through
rank four agree with the original ones. Apply
Theorem~\ref{thm:two-matroid-small} to this truncation.
\end{proof}

The six-element lemma has no rank hypothesis. The rank bound enters only
in the reduction to even conditional ground sets of size at most six.
The next unresolved step in this approach is the eight-element
antipodal-pairs inequality, which compares feasible triples and their
complementary five-sets with feasible balanced four-set partitions.

An exact obstruction prevents simply replacing the auxiliary graph by an
arbitrary three-uniform hypergraph. Partition an eight-element set as
$V=E\mathbin{\dot\cup}\{v\}\mathbin{\dot\cup}A$, with $|E|=3$ and
$|A|=4$, and take witnesses
\[
\mathcal H=\{E\}\cup\{\{v\}\cup P:P\in\binom A2\}.
\]
Every witness has a disjoint witness. Nevertheless the doubly dependent
counts are $b_3=7$ and $b_4=8$: a successful balanced partition consists
of $E\cup\{a\}$ and $\{v\}\cup(A\setminus\{a\})$ for one of the four
$a\in A$. Thus $4b_4=32<35=5b_3$. This is an obstruction to that proof
step, not a counterexample to the $2$-matroid conjecture. Indeed, its
natural complementary downset has feasible sets $\{v\}$ and $E$, but no
feasible pair $\{v,e\}$ with $e\in E$, violating
\eqref{eq:two-augmentation}. Additional augmentation structure would be
needed at the next stage.
